% --- §9-fragment-statement ---
\paragraph{Motivating conjecture (verbatim).}
Under Assumptions~\ref{ass:finite-var}--\ref{ass:generic-exposed} and \ref{ass:two-shock}, parameterize the two components of \(O(2)\) by rotation/reflection angles \(\varphi\).  Let the LP/sign feasible set on a component be \(I_{\mathrm{LP}}\), and write the report objective as
\[
\theta_w(\varphi)=a_w\cos\varphi+b_w\sin\varphi=\rho_w\cos(\varphi-\varphi_w).
\]
Let \(r_4(\varphi,\kappa;P)\) be the vector of off-diagonal fourth-cumulant residual coordinates in the structural basis and \(d_4(\varphi,\kappa;P)\) the pure-diagonal residual.  If no \(\kappa\in\mathcal K_{\delta_4}(j)\) satisfies \(d_4(\varphi,\kappa;P)=0\) at any LP support angle on the component, set the component value of \(\eta_w^*(P)\) to \(+\infty\).  Otherwise the critical leakage radius on that component is
\[
\eta_w^*(P)=
\min_{\substack{\varphi\in\arg\max_{\phi\in I_{\mathrm{LP}}}\theta_w(\phi)\\
\kappa\in\mathcal K_{\delta_4}(j),\ d_4(\varphi,\kappa;P)=0}}
\|r_4(\varphi,\kappa;P)\|_{\mathrm{off}},
\]
with the global \(O(2)\) value obtained by taking the minimum over rotation and reflection components.  In the parent corner \(\eta=0\), this reduces to the parent two-shock witness: equality holds only when an LP support angle is one of the cumulant-contact angles; otherwise the support gap is at least \(\rho_w(1-\cos d_w)\) for the distance \(d_w\) to the nearest contact angle.

\begin{theorem}[Theorem 2: two-shock angular critical-radius formula]\label{thm:two-shock-formula}
The two-shock critical leakage radius is the minimum off-diagonal fourth-cumulant residual over LP support angles satisfying pure-diagonal contact, with the global value obtained over both components.
Under Assumptions~\ref{ass:finite-var}--\ref{ass:generic-exposed} and \ref{ass:two-shock}, let \(c\in\{+,-\}\) index the rotation and reflection components of \(O(2)\), write \(Q_c(\varphi)\) for the corresponding angular parametrization, and set
\[
I_c=\{\varphi:Q_c(\varphi)\in\mathcal Q_{\mathrm{LP}}(P)\},
\qquad
M_{w,c}=\arg\max_{\phi\in I_c}\theta_w(Q_c(\phi)).
\]
On each nonempty component,
\[
\theta_w(Q_c(\varphi))
=a_{w,c}\cos\varphi+b_{w,c}\sin\varphi
=\rho_{w,c}\cos(\varphi-\varphi_{w,c}),
\]
with the zero-report case \(\rho_{w,c}=0\) interpreted by the first linear form.
Define the component critical value by
\[
\eta_{w,c}^*(P)=
\begin{cases}
\displaystyle
\min_{\substack{\varphi\in M_{w,c}\\
\kappa\in\mathcal K_{\delta_4}(j),\ d_4(Q_c(\varphi),\kappa;P)=0}}
\|r_4(Q_c(\varphi),\kappa;P)\|_{\mathrm{off}},
& \text{if the displayed feasible set is nonempty},\\[1.2em]
+\infty,&\text{otherwise}.
\end{cases}
\]
Then
\[
\eta_w^*(P)=\min_{c\in\{+,-\}}\eta_{w,c}^*(P),
\]
where a component with empty \(I_c\) contributes \(+\infty\).  In the parent corner \(\eta=0\), \(h_{\mathrm{LP}}(w;P)=h_0(w;P)\) holds if and only if some exposed LP support angle is also a cumulant-contact angle satisfying \(d_4=0\) and \(r_4=0\).  If no exposed LP support angle is a parent contact angle and \(d_w\) denotes the within-component angular distance from the exposed LP support face to the nearest parent contact angle that maximizes \(\theta_w\) among parent contacts, then
\[
h_{\mathrm{LP}}(w;P)-h_0(w;P)\ge \rho_w(1-\cos d_w).
\]
\end{theorem}

% --- §13-refutation-attempt ---
\paragraph{Refutation attempt before proof.}
The shared setup locks this conjecture as a PartialID/support-function structural time-series problem, even though the run metadata names the exact-ID cluster.  The SWIG, overlap, and do-calculus refutation primitives therefore have no treatment, outcome, or adjustment variables to instantiate.  The corresponding finite-dimensional refutation search is over the two angular components of \(O(2)\).

First, consider a candidate LP support angle \(\varphi^\star\) for which no \(\kappa\in\mathcal K_{\delta_4}(j)\) satisfies \(d_4(Q_c(\varphi^\star),\kappa;P)=0\).  This does not refute the formula: by Assumption~\ref{ass:leakage-class}, leakage tensors have zero pure diagonal coordinates, so changing \(\eta\) can never repair pure-diagonal failure.  The component critical value is correctly \(+\infty\).

Second, try to produce a mismatch by putting a small residual on the reflection component and a large residual on the rotation component.  In the two-shock parametrization this does not produce a witness.  The labelled column \(Q_c(\varphi)e_j\) traces the same unit circle on the two components, up to the fixed sign normalization in Assumption~\ref{ass:normalization}.  Since \(\theta_w\), \(A\psi_j\), and \(Z\psi_j\) depend only on that labelled column, the LP support value is represented on both components after the angular relabelling.  The reflection only signed-permutes off-diagonal fourth-cumulant coordinates involving the non-target column, preserving the pure-diagonal equations and the Euclidean off-diagonal norm.  Hence a smaller reflection residual at a non-exposed report value cannot occur inside the locked two-shock geometry.

Third, in the parent corner, setting \(\eta=0\) forces \(r_4=0\) by Assumption~\ref{ass:leakage-class}.  A searched witness with \(r_4\ne0\) is infeasible, while a witness with \(r_4=0\) and \(d_4=0\) is exactly a parent contact angle.  Thus the bounded angular search finds no counterexample to the conjectured formula, and the proof below verifies the claim.

% --- §10 Interpretation ---
\paragraph{Interpretation.}
The theorem says that the two-shock leakage diagnostic is not an optimization over the full angular contact set.  One first finds the LP/sign support angle or angles on each component, then checks whether those same angles have pure-diagonal fourth-cumulant contact.  If they do, the required leakage radius is exactly the off-diagonal fourth-cumulant residual norm at those angles; if they do not, no finite off-diagonal leakage radius can make that component attain the LP support value.  This is a computable observed-population map: on each component, the LP support angle is the projection of the report phase \(\varphi_{w,c}\) onto the closed feasible interval \(I_c\) together with any tied endpoint, and the residuals are polynomial-trigonometric functions of \((\Gamma_h,\Sigma_u,K_u)\).

The result connects the upgrade to the parent independent-shock corner.  At \(\eta=0\), equality of the LP/sign and higher-cumulant SVAR reports is not a generic consequence of sign restrictions; it requires that an exposed LP angle also diagonalize the fourth cumulant tensor with separated target cumulants.  Positive leakage widens the contact set by allowing off-diagonal co-kurtosis residuals, but it does not relax the pure-diagonal cumulant equations.

% --- §11 Proof ---
\begin{proof}
Fix \(P\), \(w\), and a component \(c\in\{+,-\}\).  By Assumptions~\ref{ass:finite-var} and \ref{ass:normalization}, \(b_j(Q_c(\varphi))=\Sigma_u^{1/2}Q_c(\varphi)e_j\) is an affine trigonometric function of \((\cos\varphi,\sin\varphi)\).  Therefore
\[
\theta_w(Q_c(\varphi))
=w^\transpose\left(e_y^\transpose\Gamma_h b_j(Q_c(\varphi))\right)_{h=0}^H
=a_{w,c}\cos\varphi+b_{w,c}\sin\varphi .
\]
If \((a_{w,c},b_{w,c})\ne0\), write \(\rho_{w,c}=(a_{w,c}^2+b_{w,c}^2)^{1/2}\) and choose \(\varphi_{w,c}\) with \(a_{w,c}=\rho_{w,c}\cos\varphi_{w,c}\) and \(b_{w,c}=\rho_{w,c}\sin\varphi_{w,c}\), giving \(\theta_w(Q_c(\varphi))=\rho_{w,c}\cos(\varphi-\varphi_{w,c})\).  If \((a_{w,c},b_{w,c})=(0,0)\), the displayed linear form remains valid and every feasible angle is an LP maximizer on that component.

By Assumption~\ref{ass:sign-zero}, \(I_c\) is the closed angular feasible set cut out by \(A\psi_j(Q_c(\varphi))\ge0\) and \(Z\psi_j(Q_c(\varphi))=0\).  Since each component is compact and the constraints are closed, every nonempty \(I_c\) has a nonempty compact maximizer set \(M_{w,c}\).  Assumption~\ref{ass:generic-exposed} excludes the only problematic case for strict statements, namely a positive-dimensional support face on which the mismatch changes while the report is constant.

Next characterize component contact.  By the formal setup and Assumptions~\ref{ass:cumulants}--\ref{ass:leakage-class}, \(Q_c(\varphi)\in\mathcal Q_\eta(P)\) if and only if \(\varphi\in I_c\) and there exists \(\kappa\in\mathcal K_{\delta_4}(j)\) such that
\[
d_4(Q_c(\varphi),\kappa;P)=0,
\qquad
\|r_4(Q_c(\varphi),\kappa;P)\|_{\mathrm{off}}\le\eta .
\]
The implication from left to right is just the coordinate definition of \(m_\eta(Q;P)=0\): the leakage tensor can absorb exactly the off-diagonal residual and, by definition of \(\mathcal L_\eta(Q)\), cannot absorb any pure-diagonal residual.  The reverse implication holds by constructing the leakage tensor whose structural-basis off-diagonal coordinates equal \(r_4(Q_c(\varphi),\kappa;P)\) and whose pure diagonal coordinates are zero.

It follows that the first leakage radius at which the component support face meets the cumulant contact set is
\[
\inf\{\eta\ge0:
\exists \varphi\in M_{w,c},\ \exists\kappa\in\mathcal K_{\delta_4}(j)
\text{ with }d_4(Q_c(\varphi),\kappa;P)=0
\text{ and }\|r_4(Q_c(\varphi),\kappa;P)\|_{\mathrm{off}}\le\eta\}.
\]
If the displayed diagonal-contact feasibility condition is empty, no finite \(\eta\) can create contact at the component support face, proving the \(+\infty\) branch.  Otherwise the infimum equals the minimum of the off-diagonal residual norm over that feasible set.  Indeed, \(M_{w,c}\) is compact, \(d_4\) and \(r_4\) vary continuously in \(\varphi\), and \(d_4(Q_c(\varphi),\kappa;P)=0\) pins \(\kappa\) to the two observed diagonal structural-basis cumulants at \(\varphi\).  Thus there is no escaping sequence in the unbounded ambient \(\kappa\)-space, and the feasible residual image is compact.

The global statement is the same calculation over \(O(2)=O_+(2)\cup O_-(2)\).  Assumption~\ref{ass:normalization} fixes the shock label after sign and permutation normalization, and Assumption~\ref{ass:two-shock} includes both components.  For a fixed labelled shock, the labelled column \(Q_c(\varphi)e_j\) has the same image on the rotation and reflection components after angular relabelling.  Hence the report \(\theta_w\) and the sign-zero feasibility constraints have the same exposed support value on the two components, while the fourth-order reflection changes only signs or orderings of off-diagonal coordinates and preserves the norm used in Assumption~\ref{ass:leakage-class}.  Taking the minimum over the two component critical values is therefore exactly the definition
\[
\eta_w^*(P)=
\inf\{\eta'\ge0:\mathcal M_w(P)\cap\mathcal Q_{\eta'}(P)\ne\emptyset\}.
\]

Finally set \(\eta=0\).  By Theorem~\ref{thm:parent-corner}, \(\mathcal Q_0(P)\subseteq\mathcal Q_{\mathrm{LP}}(P)\) and \(\mathcal Q_0(P)\) is precisely the parent independent-shock cumulant-diagonalizing sign set.  Thus \(h_{\mathrm{LP}}(w;P)=h_0(w;P)\) holds if and only if some maximizer of \(\theta_w\) over \(\mathcal Q_{\mathrm{LP}}(P)\) belongs to \(\mathcal Q_0(P)\), which is exactly the condition that an LP support angle is a parent contact angle with \(d_4=0\) and \(r_4=0\).

If no exposed LP support angle is a parent contact angle, let \(\varphi^\star\) be an exposed support angle and let \(\varphi^0\) be a parent contact angle attaining \(h_0(w;P)\) on the same connected feasible angular interval.  Write \(\alpha=|\varphi^\star-\varphi_w|\) and \(d_w=|\varphi^0-\varphi^\star|\) along the monotone branch of the interval moving away from the report phase.  Since \(\varphi^\star\) maximizes \(\rho_w\cos(\varphi-\varphi_w)\) over the feasible interval, \(0\le\alpha\le\alpha+d_w\le\pi\).  Therefore
\[
\begin{aligned}
h_{\mathrm{LP}}(w;P)-h_0(w;P)
&\ge \rho_w\{\cos\alpha-\cos(\alpha+d_w)\}\\
&=2\rho_w\sin(d_w/2)\sin(\alpha+d_w/2)\\
&\ge 2\rho_w\sin^2(d_w/2)
=\rho_w(1-\cos d_w),
\end{aligned}
\]
which is the stated parent-corner witness.  The case \(\rho_w=0\) gives the limiting zero lower bound and no strict support gap, as expected.
\end{proof}

% --- §12 Sanity check ---
\paragraph{Sanity check: parent contact at the LP support angle.}
Suppose the LP support angle on a component is unique, \(\varphi^\star=\varphi_w\in I_c\), and the observed cumulants satisfy \(d_4(Q_c(\varphi^\star),\kappa^\star;P)=0\) and \(r_4(Q_c(\varphi^\star),\kappa^\star;P)=0\) for some \(\kappa^\star\in\mathcal K_{\delta_4}(j)\).  The formula gives \(\eta_{w,c}^*(P)=0\).  By Theorem~\ref{thm:parent-corner}, the parent contact set is a subset of the LP/sign set, and because it contains the LP maximizer, \(h_0(w;P)=h_{\mathrm{LP}}(w;P)=\rho_w\).  This is exactly the parent independent-shock equality case.

If instead the nearest parent contact angle is \(\varphi^\star+d\) with \(0<d\le\pi\) and \(\varphi^\star=\varphi_w\), then \(h_0(w;P)\le\rho_w\cos d\), so the support gap is at least \(\rho_w(1-\cos d)\).  The displayed lower bound is therefore sharp in the interior-support corner.

% --- §13 Counterexample or minimality check ---
\paragraph{Minimality check: pure-diagonal contact cannot be dropped.}
The condition \(d_4(Q_c(\varphi),\kappa;P)=0\) is essential.  Consider a two-shock population and an LP support angle \(\varphi^\star\) for which every off-diagonal structural-basis cumulant coordinate vanishes, so \(r_4(Q_c(\varphi^\star),\kappa;P)=0\), but the pure diagonal coordinates are equal:
\[
K^{Q_c(\varphi^\star)}_{1111}
=K^{Q_c(\varphi^\star)}_{2222}.
\]
Then \(d_4=0\) forces \(\kappa_1=\kappa_2\), which violates the target separation requirement in \(\mathcal K_{\delta_4}(j)\) for every \(\delta_4>0\).  If one minimized only the off-diagonal residual, the reported leakage radius would be \(0\).  Under the actual assumptions, no finite off-diagonal leakage radius can repair the failed pure-diagonal separation, because Assumption~\ref{ass:leakage-class} allows leakage only in off-diagonal coordinates.  Thus the diagonal-contact and separated-cumulant clauses are not regularity decoration; they are necessary for the formula to define the same higher-cumulant SVAR contact set.

\paragraph{Substrate audit record.}
\texttt{bypass-justified}: the declaration-level audit did not surface reusable exact-identification adjustment primitives for this theorem, and the shared setup is algebraic over \(O(2)\) support functions and fourth-cumulant coordinates rather than an SCM adjustment formula.
